\problema{Largest Rectangle in Histogram}{84}{src/06-sliding-window-stack/84-largest-rectangle-histogram.cpp}{H}

Dado un arreglo \texttt{heights} con las alturas de las barras de un
histograma (cada barra de ancho \texttt{1}, todas paradas sobre la misma
línea base), encuentra el área del rectángulo más grande que se puede
formar usando una o más barras \emph{consecutivas}, con altura igual a la
barra más chica del tramo elegido.

\paragraph{La idea, con un ejemplo de la vida real.} Imagina una fila de
libros de distinto grosor parados uno junto al otro en un librero, vistos
desde arriba como barras de distinta altura. Quieres cubrir un tramo
continuo de esos libros con una tabla rectangular apoyada encima; la tabla
solo puede llegar tan alto como el libro más bajo del tramo que decidas
cubrir (si un libro es más bajo, la tabla se hunde hasta ahí). Quieres
encontrar qué tramo, y qué altura de tabla, te da el área más grande
posible.

\begin{center}
\ttfamily
\begin{tabular}{l}
alturas: 2 1 5 6 2 3\\
mejor rectángulo: barras de altura 5 y 6, área $= 5 \times 2 = 10$\\
\end{tabular}
\end{center}

\paragraph{Cómo lo resuelve el código, paso a paso.} Usamos una
\textbf{pila monótona} (monotonic stack): una pila de índices cuyas alturas
correspondientes van siempre en aumento de abajo hacia arriba.
\begin{enumerate}
  \item Recorremos las barras de izquierda a derecha, y agregamos al final
        una barra ``fantasma'' de altura \texttt{0} (un centinela) para
        forzar a que la pila se vacíe por completo al terminar.
  \item Para cada barra nueva, mientras la altura del tope de la pila sea
        mayor o igual a la altura actual, esa barra del tope ya no puede
        seguir creciendo hacia la derecha: la sacamos de la pila y
        calculamos el rectángulo más grande que la tiene a \emph{ella}
        como altura limitante.
  \item El ancho de ese rectángulo llega, por la izquierda, hasta la barra
        que quedó justo debajo en la pila (la última vez que hubo algo más
        bajo o igual), y por la derecha, hasta la posición actual
        (exclusiva).
  \item Guardamos el área más grande vista, y al final metemos la barra
        actual a la pila.
\end{enumerate}
El centinela final es lo que garantiza que ninguna barra se quede sin
evaluar: si el histograma termina con barras cada vez más altas, sin el
centinela nunca se dispararía el ``mientras'' que las saca de la pila.

\lstinputlisting{src/06-sliding-window-stack/84-largest-rectangle-histogram.cpp}

\complejidad{$O(n)$}{$O(n)$}

\paragraph{¿Qué significa esa complejidad?} Aunque hay un ciclo anidado
(el \texttt{while} dentro del \texttt{for}), cada índice entra a la pila
exactamente una vez y sale de la pila a lo más una vez, así que el trabajo
total sigue siendo proporcional a \texttt{n}, no a \texttt{n al cuadrado}.
El espacio lo usa la pila, que en el peor caso (barras siempre creciendo)
guarda todos los índices al mismo tiempo.

\paragraph{Consejos para la entrevista (Amazon).} La solución ingenua
(para cada barra, expandirse hacia los lados hasta encontrar algo más bajo)
es \texttt{O(n al cuadrado)} y suele ser el punto de partida esperado antes
de llegar a la pila monótona. Vale la pena explicar en voz alta por qué el
ancho se calcula con el índice que quedó \emph{debajo} en la pila, no con el
índice que se acaba de sacar: ese índice inferior es la barrera real por la
izquierda. El centinela de altura \texttt{0} al final es un detalle fácil
de olvidar y una fuente común de bugs (rectángulos que quedan sin evaluar).
Casos límite: arreglo vacío (área \texttt{0}), una sola barra (el área es
esa misma altura), y todas las barras con la misma altura (el rectángulo
óptimo las usa todas).
